%! Principles of Experimental Design — Unit 1, Lesson 1.6 — Group Activity
\documentclass[10pt]{article}
\usepackage{saar-article}
\usepackage{saar-boxes}
\newcommand{\TallMath}[1]{$\displaystyle #1\rule[-1.4em]{0pt}{3.2em}$}

\begin{document}

\pageheader{Unit 1, Lesson 1.6}{Group Activity}
% NAMESTRIP: no \namedateperiod here. The student writes their name once, on the
% cover this component is stapled behind. See references/conventions.md.

\vspace{0.15in}

\begin{headlinebox}{frost}
\small
\textbf{Settle It --- Harbor Point Community Pool.} Last lesson the pool's
survey found that $70\%$ of morning swimmers sleep at least $7$ hours against
$45\%$ of evening swimmers --- but the members had chosen their own swim times,
and $75\%$ of the morning swimmers turned out to be retired. This month the
director settles the question. \textbf{$90$ members volunteer.} A random number
generator assigns \textbf{$45$ to a reserved $6$~a.m.\ slot} and \textbf{$45$ to
a reserved $7$~p.m.\ slot} for four weeks; everything else about the pool stays
the same. Every volunteer keeps a sleep log. At the end, \textbf{$27$} of the
morning group averaged at least $7$ hours of sleep, and \textbf{$18$} of the
evening group did.
\end{headlinebox}

\vspace{0.03in}

\boxguard[30]
\begin{tcolorbox}[colback=white, colframe=black!40, breakable,
    title=\textbf{Tier R --- Remediate},
    fonttitle=\bfseries, arc=2mm, left=3mm, right=3mm, top=1mm, bottom=1mm]
\small
\begin{enumerate}[leftmargin=1.6em, itemsep=2pt, topsep=3pt]

  \item Name the numbers in this experiment.

        Volunteers: \blank{1.2cm} \hfill Morning group: \blank{1.2cm} \hfill
        Evening group: \blank{1.2cm} \hfill Pool members in all: \blank{1.4cm}

  \item Nobody asked the volunteers what time they would prefer.

        \textbf{(a)} Is this an experiment? \blank{2cm}

        \textbf{(b)} Who decided which group each volunteer ended up in?
        \blank{4.6cm}

  \item Of the $45$ volunteers assigned to the morning slot, $27$ averaged at
        least $7$ hours of sleep. What percent is that?

        \begin{work}
          \dfrac{27}{45} &= 0.60 \\
                         &= 60\%
        \end{work}

  \item Of the $45$ assigned to the evening slot, $18$ did. What percent is
        that?

        \begin{work}
          \dfrac{18}{45} &= 0.40 \\
                         &= 40\%
        \end{work}

  \item Which group slept more, and by how many percentage points?
        \writeline
\end{enumerate}
\end{tcolorbox}

\vspace{0.03in}

\boxguard
\begin{tcolorbox}[colback=white, colframe=black!40, breakable,
    title=\textbf{Tier A --- Approaching Proficiency},
    fonttitle=\bfseries, arc=2mm, left=3mm, right=3mm, top=1.5mm, bottom=1.5mm]
\small
\begin{enumerate}[leftmargin=1.6em, itemsep=3pt, topsep=3pt]

  \item Put the volunteers back together.

        \textbf{(a)} How many of the $90$ volunteers slept at least $7$ hours,
        and what percent is that?

        \begin{work}
          27 + 18 &= 45 \text{ volunteers} \\
          \dfrac{45}{90} &= 0.50 = 50\%
        \end{work}

        \textbf{(b)} About how many of all $1{,}500$ pool members does that
        predict?

        \begin{work}
          0.50 \times 1500 &= 750 \text{ members}
        \end{work}

  \item Name the parts of the design.

        Experimental units: \blank{4cm} \hfill Treatment: \blank{4.4cm}

        Control group: \blank{4.4cm}

  \item In a few words each, say how this design meets the four principles.

        Comparison: \blank{3.6cm} \hfill Randomization: \blank{3.6cm}

        Replication: \blank{3.6cm} \hfill Control: \blank{3.6cm}

  \item Watching gave a $25$-point gap. This experiment gives a $20$-point gap.
        Explain why the \emph{smaller} number is the one to trust.
        \writelines{2}
\end{enumerate}
\end{tcolorbox}

\vspace{0.03in}

\boxguard
\begin{tcolorbox}[colback=white, colframe=black!40, breakable,
    title=\textbf{Tier E --- Extension},
    fonttitle=\bfseries, arc=2mm, left=3mm, right=3mm, top=1.5mm, bottom=1.5mm]
\small
\begin{enumerate}[leftmargin=1.6em, itemsep=4pt, topsep=3pt]

  \item Retirement was the lurking variable last lesson, so the director blocks
        on it: $30$ of the $90$ volunteers are retired and $60$ still work full
        time. She randomizes \emph{inside} each block.

        \begin{center}
        \renewcommand{\arraystretch}{1.35}
        \begin{tabularx}{0.93\linewidth}{@{} X c c c @{}}
        \toprule
        \textbf{Block} & \textbf{In block} & \textbf{To $6$~a.m.}
          & \textbf{To $7$~p.m.} \\
        \midrule
        Retired & $30$ & \blank{1.4cm} & \blank{1.4cm} \\
        Still working full time & $60$ & \blank{1.4cm} & \blank{1.4cm} \\
        \midrule
        Total & $90$ & \blank{1.4cm} & \blank{1.4cm} \\
        \bottomrule
        \end{tabularx}
        \end{center}

        Why block on retirement instead of trusting the random assignment to
        even it out? \writeline

  \item Each change below breaks one principle. Name it, and say what could go
        wrong.

        \textbf{(a)} The director lets each volunteer pick the slot they prefer.
        \writeline

        \textbf{(b)} The morning group swims in the newly renovated lap pool;
        the evening group swims in the old one. \writeline

  \item Could this experiment be run \emph{double-blind}? Say who could be kept
        from knowing which group a volunteer was in, and who could not.
        \writelines{2}
\end{enumerate}
\end{tcolorbox}

\end{document}
